% Relative to the directory from which pdfLaTeX runs.
\makeatletter
\def\input@path{{../../../preamble/}{../../../../book2026-2027/preamble/}}
\makeatother
\input{preamble.tex}

\geometry{letterpaper,hmargin=0.75in,vmargin=0.65in}
\hypersetup{hidelinks,pdftitle={Population Parameters and Sample Statistics}}
\newif\ifprintsolutions
\printsolutionstrue % Set \printsolutionsfalse for a student-only PDF.
\ifprintsolutions
    \input{graphs/highlights.tex}
    \input{graphs/specific_highlights.tex}
\fi

% Full-width writing lines inside a task.
\newcommand{\writinglines}[1]{%
    \par\nobreak
    \begingroup
        \setlength{\parskip}{2pt}%
        \foreach \lineindex in {1,...,#1}{%
            \noindent\rule{0pt}{1.45em}\rule{\linewidth}{0.35pt}\par\nobreak
        }%
    \endgroup
}

\newcommand{\activitystage}[1]{%
    \par\addvspace{0.8\baselineskip}%
    {\large\bfseries #1\par}\nobreak\smallskip
}

\begin{document}

\begingroup
\linespread{1.08}\selectfont
\setlength{\parskip}{4pt}

{\centering
    {\Large\bfseries Population Parameters and Sample Statistics\par}
    Global Economics \textbar{} Grade 11 \textbar{} Term 2 \textbar{} Lesson 1\par
}
\medskip

\textbf{Name} \rule{0.50\linewidth}{0.4pt}\hfill
\textbf{Date} \rule{0.19\linewidth}{0.4pt}

\textbf{C1.} Distinguishes population parameters from sample statistics by interpreting and comparing the notation for population and sample means, standard deviations, and proportions in economic and financial contexts.

\activitystage{Worked example \enspace A market association's sales survey}

The \textbf{population} is all 1,200 stalls licensed for June, the group the association wants to describe. The \textbf{sample} is the 60 stalls it randomly selects. Of those 60, 18 accept digital payment.

\begin{tabularx}{\linewidth}{@{}>{\raggedright\arraybackslash}X l l@{}}
    \toprule
    \textbf{Measure} & \textbf{Population parameter} & \textbf{Sample statistic} \\
    \midrule
    Mean June sales revenue & $\mu$ & $\bar{x}=4.8$ million COP \\
    Standard deviation of June sales revenue & $\sigma$ & $s=1.2$ million COP \\
    Share that accept digital payment & $p$ & $\hat p=18/60=0.30=30\%$ \\
    \bottomrule
\end{tabularx}

The parameters describe all 1,200 stalls and their values are unknown from this survey. Among the 60 selected stalls, sales average 4.8 million COP with a standard deviation of 1.2 million COP. To get the sample's digital-payment share, divide the 18 stalls that accept digital payment by the 60 selected.

\activitystage{Guided practice}
\begin{enumerate}[leftmargin=*,series=problems]
    \item \begin{minipage}[t]{\linewidth}
        \textbf{A supermarket's checkout survey}\par
        A manager studies all 1,500 Saturday purchases at one store. Random samples A and B each contain 50 receipts from those purchases. A loyalty card was used in 20 purchases in A and 25 in B.
        \begin{enumerate}[label=(\alph*),leftmargin=*,itemsep=0.3em]
            \item Identify the population and what each sample contains.
            \writinglines{2}

            \item Label each statistic with its symbol and complete the table. Purchase values are in thousand COP. Show the proportions as percentages and give their difference in percentage points.

            \begin{tabularx}{\linewidth}{@{}>{\raggedright\arraybackslash}X p{.16\linewidth} p{.16\linewidth} p{.19\linewidth}@{}}
                \toprule
                \textbf{Measure and symbol} & \textbf{A} & \textbf{B} & \textbf{B minus A} \\
                \midrule
                Mean purchase value \rule{1.2cm}{0.35pt} & 64 & 61 & \rule{0pt}{1.6em} \\
                \midrule
                Standard deviation \rule{1.2cm}{0.35pt} & 18 & 20 & \rule{0pt}{1.6em} \\
                \midrule
                Loyalty-card share \rule{1.2cm}{0.35pt} & & & \rule{0pt}{1.6em} \\
                \bottomrule
            \end{tabularx}

            \item Use the table to describe how the three measures differ in the samples. In this survey, what do $\mu$, $\sigma$, and $p$ describe, and why do their values stay fixed across the two draws?
            \writinglines{4}
        \end{enumerate}
    \end{minipage}
\end{enumerate}
\endgroup

\clearpage
\activitystage{Independent practice}
\begin{enumerate}[leftmargin=*,resume=problems]
    \item \begin{minipage}[t]{\linewidth}
        \textbf{A credit union's loan review}\par
        On June 30, a credit union reviews its 800 active microenterprise working-capital loans. The lending committee then broadens the review to include active microenterprise equipment loans. Together, the two categories contain 1,000 active loans on that date. For each review the credit union randomly selects 50 loans from the relevant list.

        \smallskip
        \begin{tabularx}{\linewidth}{@{}>{\raggedright\arraybackslash}X >{\centering\arraybackslash}p{.19\linewidth} >{\centering\arraybackslash}p{.19\linewidth}@{}}
            \toprule
            \textbf{Measure} & \textbf{Initial sample} & \textbf{Broader sample} \\
            \midrule
            Mean outstanding balance (million COP) & 8.0 & 8.8 \\
            Standard deviation of outstanding balances (million COP) & 2.4 & 2.9 \\
            Loans with an overdue installment on June 30 & 10 & 8 \\
            \bottomrule
        \end{tabularx}

        \begin{enumerate}[label=(\alph*),leftmargin=*,itemsep=0.8em]
            \item Identify the population and sample in each review. Calculate the change in population size.
            \writinglines{4}

            \item Give the symbol and meaning of each population parameter in the two reviews. Pair each parameter with its sample statistic. Calculate both sample proportions and interpret the three statistics in each sample.
            \writinglines{8}

            \item Calculate each difference (broader sample minus initial sample), with units. Express the difference in overdue shares in percentage points. Explain what the data let you conclude about changes in the population parameters.
            \writinglines{6}
        \end{enumerate}
    \end{minipage}
\end{enumerate}

\ifprintsolutions
    \clearpage
    \setlength{\parskip}{4pt}
    \section*{Solutions}
    \PracticeHighlightLegend{G11-T2-C1}

    \begin{enumerate}[leftmargin=*,itemsep=0.8em]
        \item \textbf{A supermarket's checkout survey}
        \begin{enumerate}[label=(\alph*),leftmargin=*,itemsep=0.3em]
            \item The population is \HighlightStroke{G11T2C1Population}{all 1,500 Saturday purchases} at the store. Samples A and B are two draws of \HighlightStroke{G11T2C1Sample}{50 Saturday purchases}, represented by the receipts.

            \item Each difference is B minus A. The mean and standard deviation are in thousand COP. The share difference is in percentage points (pp).

            \begin{tabularx}{\linewidth}{@{}>{\raggedright\arraybackslash}X l l l@{}}
                \toprule
                \textbf{Statistic} & \textbf{A} & \textbf{B} & \textbf{B minus A} \\
                \midrule
                Mean $\HighlightStroke{G11T2C1SampleMean}{\bar{x}}$
                  & $\HighlightStroke{G11T2C1SampleMean}{64}$
                  & $\HighlightStroke{G11T2C1SampleMean}{61}$ & $61-64=-3$ \\
                Standard deviation $\HighlightStroke{G11T2C1SampleSD}{s}$
                  & $\HighlightStroke{G11T2C1SampleSD}{18}$
                  & $\HighlightStroke{G11T2C1SampleSD}{20}$ & $20-18=+2$ \\
                Loyalty-card share $\HighlightStroke{G11T2C1SampleProportion}{\hat p}$
                  & $20/50=\HighlightStroke{G11T2C1SampleProportion}{40\%}$
                  & $25/50=\HighlightStroke{G11T2C1SampleProportion}{50\%}$ & $50-40=+10$ pp \\
                \bottomrule
            \end{tabularx}

            \item Sample B's mean purchase value is 3 thousand COP lower, its standard deviation is 2 thousand COP higher, and its loyalty-card share is 10 percentage points higher. Across all 1,500 purchases, $\HighlightStroke{G11T2C1PopulationMean}{\mu}$ is the mean purchase value, $\HighlightStroke{G11T2C1PopulationSD}{\sigma}$ is the standard deviation of purchase values, and $\HighlightStroke{G11T2C1PopulationProportion}{p}$ is the loyalty-card share. Each parameter is fixed because both samples come from the same set of purchases. The parameter values themselves are unknown.
        \end{enumerate}

        \item \textbf{A credit union's loan review}
        \begin{enumerate}[label=(\alph*),leftmargin=*,itemsep=0.3em]
            \item On June 30 the initial population is \HighlightStroke{G11T2C1Population}{all 800} active microenterprise working-capital loans. The broader population is \HighlightStroke{G11T2C1Population}{all 1,000} active loans across both categories. Each sample consists of \HighlightStroke{G11T2C1Sample}{50 loans} drawn from its population. The expanded population adds $1{,}000-800=200$ loans.

            \item Each parameter describes all loans in the relevant population (800 or 1,000). Each statistic describes the 50 selected loans in its column. The two balance measures are in million COP.

            \begin{tabularx}{\linewidth}{@{}>{\raggedright\arraybackslash}X l l@{}}
                \toprule
                \textbf{Measure and paired symbols} & \textbf{Initial sample} & \textbf{Broader sample} \\
                \midrule
                Mean outstanding balance
                  $\HighlightStroke{G11T2C1PopulationMean}{\mu}$ and $\HighlightStroke{G11T2C1SampleMean}{\bar{x}}$
                  & $\HighlightStroke{G11T2C1SampleMean}{8.0}$
                  & $\HighlightStroke{G11T2C1SampleMean}{8.8}$ \\
                Standard deviation of outstanding balances
                  $\HighlightStroke{G11T2C1PopulationSD}{\sigma}$ and $\HighlightStroke{G11T2C1SampleSD}{s}$
                  & $\HighlightStroke{G11T2C1SampleSD}{2.4}$
                  & $\HighlightStroke{G11T2C1SampleSD}{2.9}$ \\
                Share with an overdue installment
                  $\HighlightStroke{G11T2C1PopulationProportion}{p}$ and $\HighlightStroke{G11T2C1SampleProportion}{\hat p}$
                  & $10/50=\HighlightStroke{G11T2C1SampleProportion}{20\%}$
                  & $8/50=\HighlightStroke{G11T2C1SampleProportion}{16\%}$ \\
                \bottomrule
            \end{tabularx}

            \item The broader sample has a higher mean balance ($8.8-8.0=+0.8$ million COP), a higher standard deviation ($2.9-2.4=+0.5$ million COP), and a lower overdue share ($16\%-20\%=-4$ percentage points). The parameters describe a broader population. Whether or by how much their values differ is unknown from the samples alone.
        \end{enumerate}
    \end{enumerate}
\fi

\end{document}
